% Section fragment for the Next manuscript.
% Requires amsmath, amssymb, and algorithmic in the main preamble.
% The pseudocode blocks use IEEEtran figure floats, not algorithm floats.

\section{协议设计}
\label{sec:protocol}

Next 认证固定的支付授权与输出，并允许它们的来源保持未解决状态。公开执行以原子方式消费缺失的已认证输入并记录发行方的直接义务，同时公布来源交易的证书，使原签名者能够重新提交该来源交易。截止期时仍缺失的来源经由无需门限适配的公开赔付决定而关闭，表示命令随后在既有承诺之下替换资金引用。

\subsection{身份、认证与交付}
\label{sec:protocol-certification}

令 $K_T$ 为编码后的支付体与输入声明，$\mathbf{C}_T$ 为输入承诺，$\Sigma_T$ 为输出摘要：
\begin{equation}
  \begin{aligned}
    K_T &= \operatorname{Encode}(T.\mathsf{Body},T.\mathsf{Claims}),\\
    \mathsf{id}(T) &= H(\texttt{TX\_V3},K_T,\mathbf{C}_T),\\
    c_T &= H(\texttt{DIRECT\_CERT\_V3},\operatorname{Encode}(\Sigma_T)).
  \end{aligned}
  \label{eq:protocol-identities}
\end{equation}
所有者签名认证 $\mathsf{id}(T)$ 且不参与其计算，每个输出标识符绑定该身份、网络和输出位置。资金引用与开口单独编码，因此受限替换不改变已授权的声明与输出。初始提交要求每个被声明的输入具有 \texttt{OriginalFunding} 及有效开口，而后续表示保持支付体、声明、承诺和所有者签名不变。因此验证缓存以完整命令字节为键。

在同一 Network 与 Subject 之下，诚实钱包对每笔新支付（跨组织与实例）使用新的支付 nonce；重试则以其原授权重新提交同一笔支付。钱包不得更改 nonce 而复用旧证书，共享同一 intent 的两个事实也不会被合并为一项担保。实验驱动程序从输入标识符导出 nonce，这并不是通用的多设备 nonce 分配服务。

成员检查所有权、路由、直接输入证书、价值守恒和准入向量。一个本地事务重新检查输入锁、最终费用输入以及指定 Worker 的资源切片，并连同其精确的扣减一起记录批准；该批准保留请求的交易、准入向量和直接输入证书（第~\ref{sec:protocol-resubmission}~节）。即使每个本金输入都是最终的，CAL 预留也等于包含找零的 $a_T$，因为每个已认证输出都可能在该支付执行之前流转。成员在该事务提交之后签名，重试复用未被作废的批准，而释放资源额度则保留消费锁。

收集者向四名成员发送一个规范请求，在验证了针对同一 $c_T$ 的三个签名后返回；不同的有效子集给出同一事实。接收方对照输出体、索引、所有者和可信配置检查 TXCer，原子地记录该币，并可立即构造后继请求，其中携带其直接输入证书以代替祖先链。

中继材料在后台保存，有界的早期任务使 \texttt{INSTALL} 与公开提交能够相互独立地进行。\texttt{INSTALL} 在成员处存储完整材料并在本地记录已认证的消费；它既不创建新的批准，也不产生预算扣减。公开提交单独进行。成员相对于其第一次 \texttt{INSTALL} 来安排备份提交，重复的 \texttt{INSTALL} 不能推迟它。HTTP 接受使发件箱条目保持待处理状态，直到经认证的公开完成将其移除。备份提交需要保留的完整提交；第~\ref{sec:protocol-resubmission}~节增加了一条基于原签名者批准的路径。

公开的 \texttt{DirectSubmission} 携带交易、准入向量、组织 \texttt{SpendQC} 和直接输入证书。委员会重建 $\Sigma_T$ 并验证同样的授权；省略单独的新输出 TXCer 封装只改变信封而保持该要求。公开验证将每份输入证书绑定到其网络、发行方、事实和输出索引。每个认证输入必须恰有一个证书与索引条目；重复的 OutputID 和多余条目被拒绝，同一父证书的不同输出可以分别使用。覆盖按父事实登记一次。准入向量含一个 CAL 条目，且绑定到发行方。

\subsection{带直接义务的原子执行}
\label{sec:protocol-execution}

公开执行将整份证书范围的覆盖与实际的缺失输入义务区分开。在父事实 $c$ 第一次被缺失使用时，\texttt{register} 预留其全部已认证的 CAL 金额并为每个输出创建一个 Promise。覆盖记录维持
\begin{equation}
  \mathsf{Original}_c=
  \mathsf{Discharged}_c+\mathsf{Paid}_c+\mathsf{Remaining}_c,
  \label{eq:protocol-coverage}
\end{equation}
其中 $\mathsf{Paid}_c$ 是毛额累计计数器，$\mathsf{Rec}_c\le\mathsf{Paid}_c$ 统计由来源执行偿还的部分。Grant 用量满足 $\mathsf{Reserved}_g+\mathsf{Spent}_g\le B_g$，其中 $\mathsf{Spent}_g$ 对 $g$ 之下各证书的 $\mathsf{Paid}_c-\mathsf{Rec}_c$ 求和。登记提高公开的 grant 用量而不触动备付账户；一次生效的赔付（第~\ref{sec:protocol-decision}~节）扣减该账户并提高 $\mathsf{Spent}_g$，偿还则贷记该账户并将 $\mathsf{Spent}_g$ 降低同样的数额。只有实际被消费的缺失输出才会产生直接义务并提高公开缺口，因此消费一份 40/60 证书中缺失的一个 40-CAL 输出，会预留 100~CAL、产生一个 40-CAL 的缺口，并使未使用的 Promise 保持被覆盖。

图~\ref{fig:protocol-payment} 给出公开转换。其前置状态包含同一区块中较早的成功命令；静态授权可被缓存，而消费、余额、grant 和义务则对照当前状态评估。所有变更都发生在私有 overlay 中，失败将丢弃整个转换。

intent 检查属于此评估：若 $T$ 的 intent 已被绑定到不同的 TxID，则该转换返回冲突。该绑定在成功执行一笔支付的同一原子转换中写入，且没有任何生产路径删除或重新绑定它，因此重新提交相同字节不能消除该冲突。

\begin{figure}[t]
  \small
  \begin{algorithmic}[1]
    \REQUIRE 已验证的提交 $T$，当前状态 $S$，公开时间
    \STATE 若 $T$ 已 settled，则返回无新效果
    \STATE $O\gets\operatorname{Overlay}(S)$
    \STATE 检查 intent、grant 和 work；在 $O$ 中为费用托管注资
    \FOR{每个 CAL 输入 $(x,i)$}
      \STATE 要求 $(x,i)$ 未被消费
      \IF{存在 $(x,i)$ 的最终创建}
        \STATE 对照声明检查其区块体与证据
      \ELSE
        \STATE 要求 $i=0$ 且存在匹配的直接 TXCer $c$
        \STATE 若尚未见过，则登记 $c$ 的全部覆盖
        \STATE 创建唯一的 $\mathcal{O}_x$；增加缺口与 Pending
      \ENDIF
      \STATE 在 $O$ 中记录 $(x,i)$ 的消费
    \ENDFOR
    \STATE 要求 CAL 输入总额与输出总额相等
    \STATE 记录 $T$ 为 settled；推进其费用阶段
    \STATE 处理每个输出：履行 Open 义务；对 Repaired 义务向发行方偿还并将其标记为 Recovered；否则履行任何已登记的 Promise。除非已偿还，否则创建最终的实例~0
    \STATE 将所有变更与结果作为一个转换返回
  \end{algorithmic}
  \caption{公开支付执行。任何检查失败都会拒绝所有临时变更。输出处理遵循第~\ref{sec:protocol-closure}~节。}
  \label{fig:protocol-payment}
\end{figure}

义务绑定被消费的输出、父证书事实、发行方、金额、消费交易以及历史输入位置，其截止期以公开区块时间为锚。消费支付记录其 Pending 输入义务与费用进度。用户支付的费用消费经认证的最终 FUEL，将多余的输入价值作为找零返还，并托管已签名的最大值，该最大值涵盖登记、结算、闭合以及每个已认证输入的一次赔付费用；唯一的阶段与义务标识符防止重复收费，未使用的托管在输入义务闭合时返还给付款方。即使部分输入仍是义务，子输出也是最终的，因此后续交易可依据普通的最终输入检查消费它们。

\subsection{来源闭合、偿还与累计对账}
\label{sec:protocol-closure}

来源交易通过普通的支付执行到达，包括验证并消费其自身的输入，其输出根据各自的义务在该转换内被处理。到达将 Open 义务关闭为 Fulfilled：它履行 Promise、解除覆盖并降低消费方的 Pending 计数。没有任何支付消费的输出会履行任何已登记的 Promise 并成为最终的实例~0。Repaired 义务意味着由备付扣减为消费方提供了资金，因此到达时偿还该扣减，而不是创建第二个输出：金额返回到付款的发行方账户，义务变为 Recovered，偿还计数器与 grant 的 Spent 同时变化，而赔付决定保持原有记录。已被消费的实例~0 在每个分支中都保持已消费；被偿还的输出没有最终创建。来源交易自身的输入可能在同一原子转换内针对其直接发行方产生新的义务。

成员根据经认证的连续区块和结果重建进度，且只有已存在的本地批准才提供待恢复的扣减。对其 CAL 扣减，成员 $m$ 计算
\begin{equation}
  \begin{aligned}
    t_{m,T}&=
    \begin{cases}
      a_T-\bigl(\widehat{p}_{m,T}-\widehat{\mathsf{Rec}}_{m,T}\bigr),
        &\text{若 settled},\\
      0,&\text{否则},
    \end{cases}\\
    \delta_{m,T}&=t_{m,T}-\mathsf{Applied}_{m,T},
  \end{aligned}
  \label{eq:protocol-release}
\end{equation}
其中 ``settled'' 表示该成员已观察到自己的支付成功执行，$\widehat{p}_{m,T}$ 与 $\widehat{\mathsf{Rec}}_{m,T}$ 是为该批准的输出在本地记录的赔付与偿还。在观察到其支付的执行时，成员检查经认证结果中的偿还额覆盖 $\widehat{p}_{m,T}$，并令 $\widehat{\mathsf{Rec}}_{m,T}=\widehat{p}_{m,T}$，因为该次执行偿还了 $T$ 的每个已赔付输出；在跟随赔付之后才批准 $T$ 的成员则未记录任何赔付。在检查边界与单调性之后，它将 $\delta_{m,T}$ 在原 Worker 中从 Reserved 移至 Available 并更新 Applied。Pending 输入不影响 CAL 释放，因为这些责任由其发行方承担，而 work 资源则等待费用闭合。未执行的已认证来源在每个批准成员处保持预留其全部 CAL 金额（含找零）；仅有赔付不会释放它。已花费的 FUEL 不会作为未花费的费用被恢复。重复的区块与迟到的材料不能重复贷记。钱包将被偿还的输出记录为终态观察，其没有证书、最终创建或余额。

经授权的备付金增加从外部的、未受保护的 CAL 账户扣款，贷记备付金，并以原子方式提高其 grant。每个成员加上新旧 $\lfloor2B_g/3\rfloor$ 份额之差，同时 Reserved 及其批准保持不变，汇总的支持检查覆盖共享同一账户的所有 grant。

\noindent\textbf{被取代的部分批准。}
公式~\eqref{eq:protocol-release} 规范正常的结算与偿还。另有一个独立的本地终态转换，适用于以下情形：同一组织配置下经认证且成功应用的公开支付，消费了成员曾为另一个事实锁定的本金输入实例或用户支付的费用输入实例。在替换本地花费候选之前，跟随者重新加载不可变的批准，并检查确切的输出实例、不同的事实与交易，以及不存在已观察到的执行、已安装或已排队的证书材料、待处理义务、赔付、恢复或任何费用效果。

随后，成员对每个已记录的本地扣减 $d$，在原资源切片中恰好将 $d.\mathsf{Cap}-\mathsf{Applied}$ 从 Reserved 归还到 Available，并将该批准记录为被第一个成功的冲突事实及其高度所作废。它不设置 \textsf{Settled}，不关闭费用托管，不修改批准，不清除其他输入锁，也不改变公开状态。累计的 $\mathsf{Applied}$ 向量使该转换在同一批准的多个输入发生冲突或区块被重放时保持幂等。重试与迟到的 \texttt{INSTALL} 被拒绝。将已作废事实作为所谓已认证来源而暴露，被视为会计矛盾，并在所述故障模型下使跟随者停止。超时、仅 \texttt{INSTALL}，或不相交输入之间的全局 Intent 冲突，均不授权此释放。

\subsection{基于公开证书的来源重新提交}
\label{sec:protocol-resubmission}

执行某笔支付的区块携带认证其已认证输入及其创建支付的法定人数证书。每个成员在其经认证的区块事务中跟随该区块。对于所携带的、属于其自身组织的每份证书，它加载以该证书的事实为索引的批准，除非它已观察到该来源的公开执行，或已持有该事实的发件箱条目，并检查所存储的交易与准入向量能复现该事实。随后它组装所存储的交易、所存储的直接输入证书和暴露的 QC，并应用普通的支付验证。成功时，它在同一个原子本地更新中写入一个发件箱条目，并按成员索引错开首次尝试；既有的中继在提交后提交该条目。其他组织的支付触发同一路径，而未通过这些检查的已存储记录属于损坏状态，会使跟随者停止。

该作业重新提交该组织所认证的支付，因此不需要新的所有者或成员签名，不创建批准，保持输入锁不变，也不预留额度。已存在的发件箱条目保留其重试时间，来源交易的执行会移除该条目。该路径需要持有该批准、跟随链并能到达委员会的原签名者，也需要不存在阻止该来源执行的全局 intent 冲突；截止期时仍缺失的来源交由赔付处理。它不提供固定时间保证。第~\ref{sec:security-decision}~节表明，只要义务已公开，至少有两名诚实签名者持有该材料。

\subsection{赔付决定}
\label{sec:protocol-decision}

截止期之后，Open 义务通过公开的 \texttt{CompensationDecision} $\delta=(\mathit{net},x,h,j,k)$ 关闭：它由网络、被消费的输出，以及消费支付的高度、交易索引和输入索引构成。定长的规范编码既不含提交者签名，也不含适配数据。权限来自公开义务，且已提交状态固定目标坐标，因此对于给定的义务，每个提交者构成相同的字节。委员会 worker 扫描到期索引并提交到期的决定，而无需请求适配份额。

图~\ref{fig:protocol-decision} 给出该转换。执行器在读取任何历史区块之前，先将请求的网络、位置和输入槽与义务比较，因此不匹配的请求无需加载历史即被拒绝。随后它读取位于 $(h,j)$ 的规范消费支付，检查其身份、事实、输入和被声明的金额，并要求 $x$ 具有 \texttt{OriginalFunding}、有效的开口以及同宽度的备付替换；它不计算替换开口，也不联系适配器。一个 overlay 承载所有变更：备付扣减、Repaired 状态、Paid、Spent、缺口、消费方的 Pending 计数、赔付费用（从消费方托管到奖励），以及在没有剩余输入义务时的消费方费用闭合与退款，连同不可变的决定记录 $\mathcal D_x$（义务快照与备付扣减标识）和待表示条目。

\begin{figure}[t]
  \small
  \begin{algorithmic}[1]
    \REQUIRE 决定 $\delta=(\mathit{net},x,h,j,k)$，状态 $S$，执行高度与区块时间
    \STATE 要求 $h$ 早于执行高度
    \STATE 若 $\mathcal D_x$ 存在，则当其等于 $\delta$ 时返回无效果；否则拒绝
    \STATE 要求 $\mathcal O_x$ 为 Open，消费方位置为 $(h,j,k)$，网络为 $\mathit{net}$
    \STATE 要求 $\mathcal O_x$ 的截止期已被区块时间到达
    \STATE 加载位于 $(h,j)$ 的规范消费支付；检查其身份、事实、输入和被声明的金额；要求 $x$ 具有 \texttt{OriginalFunding}、有效开口以及同宽度的备付引用
    \STATE $O\gets\operatorname{Overlay}(S)$
    \STATE 在 $O$ 中：扣减发行方备付金；将 $\mathcal O_x$ 设为 Repaired；将其 Promise 转为 Paid；更新 Spent、缺口以及消费方的 Pending 计数和费用阶段
    \STATE 在 $O$ 中：记录 $\mathcal D_x$ 和待表示条目
    \STATE 返回变更与赔付结果
  \end{algorithmic}
  \caption{赔付决定。该转换不读取也不请求适配材料。任何检查失败都会拒绝所有临时变更。}
  \label{fig:protocol-decision}
\end{figure}

结果携带单一的逐输出效果（它指明父事实与消费方事实、金额、输入位置和扣减标识）以及任何费用输出。重复的相同决定返回一个变更集为空的未应用结果（在来源偿还之后亦然），与记录不同的决定则被拒绝；一旦来源已执行，针对其义务的新决定不再提交。成员连同其游标一起应用经认证结果中的效果，并定位原始批准以进行 Paid 与 Pending 更新，钱包对每个费用输出只应用一次。赔付不向接收方或已执行的后继交易计入任何本金。

\subsection{表示命令与批次}
\label{sec:protocol-batch}

一个 \texttt{RepairBatch} 针对一个历史区块，按规范顺序含 1--32 个条目，每个条目指明一个输出、一个交易索引、一个输入索引和一个新的输入开口；输出与槽位唯一。该批次固定基础修订以及先前与替换后的规范区块体摘要，携带 part 开口，并将逻辑修订从 $\mathsf{Base}$ 推进到 $\mathsf{Base}+1$。单个 \texttt{RepairInput} 是其单条目形式，并携带被替换的交易字节。两种命令都只接受具有已提交赔付决定的输出，从决定记录中获取金额、发行方和替换引用，并且不写入会计状态。

公开执行对照其当前状态重复所有检查。基础已变化、某条目没有决定，或某槽位不再持有 \texttt{OriginalFunding}，都会拒绝整个批次，且确切的成功重放在基础验证之前即被识别。由于每条命令都要求其槽位中为 \texttt{OriginalFunding}，不同的批次或单条目入口不能对同一槽位替换两次，命令与修订标识保持特定于所选的打包方式。经认证的结果指明该命令及其是否被应用；跟随者将其映射为空的经济效果，而赔付的经济效果仅来自决定结果。偿还属于来源交易自身的执行结果。

\subsection{保持身份的安装与重放}
\label{sec:protocol-installation}

四种身份被区分开。支付身份 $\mathsf{id}(T)$ 固定授权与承诺，且所有者签名绑定它；它不是共识交易哈希。格式良好的 \texttt{UTXO4CH} 信封的共识 $\mathit{Tx.Hash}$ 是对其头部与固定字段的 SHA-256，不检查任何支付有效性。$\mathit{DataHash}$ 是共识交易哈希之上的 Merkle 根，而包含性证明不授予修订权限。完整 BlockID $(H_b,P_b)$ 保留区块哈希与原始 part-set 头，执行与安装都比较这两者。Part 承诺是绑定到链、密钥、高度和 part 索引的变色龙承诺，被更改的 part 携带指向其原始承诺的有效开口。原始共识 part 携带 revision~0 与 opening~1；原始验证要求正确的高度与这些值，且应用正常验证原始支付。被修订的 part 不能冒充新的原始执行，只有已提交的 Task 才授权改写。重放使用保留的原始命令，且从不将被修订的区块体作为新支付执行。

公开的修订与 Task 授权一个确切的新区块体及其 part 证据，安装器在获取物理修订锁之前，从一个已提交的读视图中复制它们。对于目标修订 $r$，较低的物理版本被验证并安装，相等的版本要求材料相同，较高的版本则不经回滚而覆盖旧任务。原始区块被保留，part 与本地版本元数据被原子地写入。安装可以跳过中间的物理版本，并保留每条经济命令。队列进度保持全局且连续：对于条目 $A_1,B_1,A_2$，在处理 $A_1$ 时安装 $A_2$ 会使 $B_1$ 保持未完成。逻辑执行读取已提交的 Canonical 修订，因此每种安装调度都留下相同的经济历史。当某个修订的区块体已在某副本上被物理安装，该修订即在该副本处被\emph{物化}。

重放按已提交的顺序执行保留的原始区块，包括支付、赔付决定和表示命令，并重建该序列的应用哈希。既有的后继交易保持其授权、输出引用和区块身份。仅赔付决定就以仅追加的公开记录在经济上关闭义务。表示增加了一个接口：在消费方的原始坐标 $(h,j,k)$ 处，已存储的区块指明指定输入的备付扣减，并仍能对照原始 BlockID 验证。读取者仍读取永久决定与义务状态以获得经济状态，被修订的字段并不取代它们。批处理在一个区块内共享表示工作，直接安装合并历史写入，二者都不为正常认证增加一轮交互。第~\ref{sec:eval-reader}~节报告了该接口的功能区别，以及相对于优化后决定索引的读取、存储与维护权衡。

\subsection{在原始坐标处读取资金引用}
\label{sec:protocol-reader}

位于 $(h,j,k)$ 的读取返回 $(H,h,j,k,F,r,d,s_H)$，由现有函数在 $V_H$ 上的一个读事务内组装而成。其中 $F$ 是资金表示，$r$ 是其修订，$d$ 是决定或缺失标记，$s_H$ 是义务状态。该逻辑元组描述本地组合，而不是已部署的网络接口。索引 $j$ 对原始交易列表计数，不会因过滤维护命令或失败交易而重新编号。逻辑区块体来自 $V_H$ 中已提交的修订，仅当不存在修订时才回退到原始区块，因此物理安装进度，或重放期间已持有较晚版本的区块存储，都不改变答案。存在但无法解码的修订是错误，绝不回退到原始区块。缺失的历史被报告为不可用，而不被解读为不存在赔付；该契约假设相关历史未被裁剪。在完整、已验证、未裁剪的 $V_H$ 中，缺失决定意味着截至 $H$ 不存在决定。

\noindent\textbf{四个状态维度。}
这些维度描述不同的对象并各自独立推进。

\par\smallskip
{\small
\setlength{\tabcolsep}{3pt}
\noindent
\begin{tabular}{@{}p{0.24\columnwidth}p{0.70\columnwidth}@{}}
\toprule
维度 & 含义 \\
\midrule
公开支付 &
\texttt{Settled} 记录成功执行。新创建的输出即使在 \texttt{Pending} 非零时也是最终的；具有 Repaired 义务的输出向备付金偿还，而不是被重新创建。 \\

直接义务 &
Open$\to$Fulfilled，或 Open$\to$Repaired$\to$Recovered。
Repaired 表示赔付；Recovered 表示随后的偿还。
赔付决定持续存在。 \\

逻辑表示 &
$V_H$ 中的修订决定该槽位。
\texttt{OriginalFunding} 并不意味着不存在赔付。
\texttt{ReserveFunding} 记录一次历史垫付，并在偿还之后仍然保留。 \\

物理安装 &
$\rho_p$ 记录已安装的区块存储版本。
安装进度不决定 $H$ 处的逻辑答案，也不改变经济状态。 \\
\bottomrule
\end{tabular}
\par}

因此，读取者将永久决定与义务状态连同资金表示一并查阅。待表示条目是可移除的工作状态，而不是赔付已发生的持久证据。这些表示或安装变更都不会向消费方计入第二笔本金。

\noindent\textbf{用途：以区块记录为单位的归档复核。}
考虑一种本地归档复核：它保留每个区块的原始提交证据与交易坐标。赔付发生后，其输入复核需要将一次备付垫付关联到消费了缺失输出的支付。索引或物化资金视图可以在不改变原始正文的情况下，以附注提供这一关联。C3 提供另一种记录格式：指定输入字段本身指明历史备付扣减，授权区块体仍与保存的提交证据相容。因此，复核者可替换正文表示，同时保留原区块记录及其已有提交证据。

消费者从 $V_H$ 读取授权正文与确切的 Revision 和 Task，使用该修订的开口重建区块与 part，并通过 Next 的验证接口，对照保存的提交证据检查完整 BlockID。投票同时签署区块哈希与 part-set header。由于资金槽位位于共识交易哈希覆盖的固定字段之外，改变该槽位仍保持头部哈希；若未适配 part，part-set 检查则拒绝修改后的正文。消费者随后提取 $(h,j,k)$ 处的备付扣减引用，并将其与永久决定和 $H$ 处的义务状态关联。旧提交证据确定与已提交身份的相容性；成功执行的决定与表示命令连同 $V_H$ 中记录的确切 Task 授权新的字节。

原型执行了这一流程，包括迟到偿还。表~\ref{tab:c3-function} 将这种正文记录格式与经济答案区分开来；索引与普通视图同样提供经济答案，并保留原身份及后继引用。消费者使用 Next 接口，并非未经修改的 CometBFT 客户端。“可选”表示后续表示服务停止时经济闭合仍可推进；当前构建依然使用 CH 交易与 part 编码。

\noindent\textbf{密码学实例化。}
该实现使用 RSA 变色龙哈希，模数为 2048 位，$e=65537$，并使用带上下文绑定的 SHAKE256 全域哈希映射到 $\mathbb Z_N^*$。一个可信 dealer 使用 CIRCL~v1.6.5~\cite{circl} 生成三选四的门限 RSA 密钥，遵循 Shoup 的门限构造~\cite{shoup}。公开的适配比值不传递编辑权限；传递编辑权限的是已提交的决定、确切的 Task 以及完整身份守卫。补充材料规定了编码、参数假设和原始群适配接口。
